\providecommand{\chLogicExample}[3]{\chThreeExample{#2}{#3}\label{ex:#1}}
\subsection{ধরন অনুযায়ী বুলিয়ান রাশির সমাধান}
প্রথমে রাশির পূরকের সীমা এবং বন্ধনী দেখো। তারপর কোন সূত্রে পদ কমবে তা নির্বাচন করো। নিচে প্রতিটি ধরনের দুইটি সমাধান রয়েছে; সমাধানের পরে সূত্রের কারণ বলা হয়েছে। বুলিয়ানে $+$ হলো OR, পাশাপাশি লেখা চলক বা $\cdot$ হলো AND।
\Needspace{12\baselineskip}
\subsubsection*{ধরন ১: অপারেশনের ক্রম ও মান বসানো}
আগে বন্ধনী ও পূরক, তারপর AND, সবশেষে OR করো। সাধারণ যোগের carry এখানে নেই।
\par\smallskip\noindent\begin{minipage}{\linewidth}
\chLogicExample{BA01.1}{১.১}{একই ১-এর দুই রকম যোগ}
$A=B=1,C=0$ হলে $F=A+B\overline C$ এবং বাইনারি যোগ $A+B$-এর ফল তুলনা করো.
\chThreeAnswer $\overline C=1$, তাই $B\overline C=1$ এবং বুলিয়ান $F=1+1=1$। কিন্তু বাইনারি গাণিতিক যোগে $1_2+1_2=10_2$। অপারেশন বোঝার আগে সংখ্যার চিহ্ন দেখে যোগ করা যাবে না।

\end{minipage}\par\smallskip
\par\smallskip\noindent\begin{minipage}{\linewidth}
\chLogicExample{BA01.2}{১.২}{আংশিক পূরকের মান}
$A=0,B=1,C=1$ হলে $F=(A+\overline B)C+\overline{AC}$ নির্ণয় করো.
\chThreeAnswer $\overline B=0$, $(A+\overline B)C=0$ এবং $AC=0$ থেকে $\overline{AC}=1$। তাই $F=0+1=1$। শেষ পূরকটি পুরো $AC$-এর।

\end{minipage}\par\smallskip
\Needspace{12\baselineskip}
\subsubsection*{ধরন ২: দ্বৈত রূপ নির্ণয়}
AND ও OR পরস্পর বদলাও; ০ ও ১ বদলাও। চলক এবং চলকের পূরক যেমন আছে তেমন থাকবে। দ্বৈত রূপকে মূল রাশির পূরক ভাববে না।
\par\smallskip\noindent\begin{minipage}{\linewidth}
\chLogicExample{BA02.1}{২.১}{বন্ধনীসহ দ্বৈত রূপ}
$F=(A+\overline B)(C+0)$-এর দ্বৈত রূপ লেখো.
\chThreeAnswer $F^d=A\overline B+C\cdot1=A\overline B+C$। $\overline B$ একই থাকে। মূল রাশিটি $C(A+\overline B)$, তাই দ্বৈত রূপের সঙ্গে সব input-এ সমান হওয়ার প্রয়োজন নেই।

\end{minipage}\par\smallskip
\par\smallskip\noindent\begin{minipage}{\linewidth}
\chLogicExample{BA02.2}{২.২}{ধ্রুবকসহ সূত্রের দ্বৈত}
$A\cdot1+AB=A$ সূত্রের দ্বৈত লিখে যাচাই করো.
\chThreeAnswer দ্বৈত হলো $(A+0)(A+B)=A$। বামপক্ষ $A(A+B)=A+AB=A$। একটি সত্য সূত্রের দ্বৈতও সত্য; চলকের নাম বদলানোর প্রয়োজন নেই।

\end{minipage}\par\smallskip
\Needspace{12\baselineskip}
\subsubsection*{ধরন ৩: সম্পূর্ণ রাশির পূরক}
ডি-মরগ্যানে OR-এর পূরক AND হয়, AND-এর পূরক OR হয়। প্রতিটি পদের পূরক নিতে হবে; দ্বৈত পূরক উঠে যায়।
\par\smallskip\noindent\begin{minipage}{\linewidth}
\chLogicExample{BA03.1}{৩.১}{যোগফলের পূরক}
$Y=\overline AB+BC$ হলে $\overline Y$ নির্ণয় করো.
\chThreeAnswer $\overline Y=\overline{\overline AB}\,\overline{BC}=(A+\overline B)(\overline B+\overline C)$। একই চলক বারবার থাকলেও প্রতিটি পদের পূরক আলাদাভাবে নেওয়া হয়েছে।

\end{minipage}\par\smallskip
\par\smallskip\noindent\begin{minipage}{\linewidth}
\chLogicExample{BA03.2}{৩.২}{গুণফলের পূরক}
$Y=(A+\overline C)(\overline B+D)$ হলে $\overline Y$ বের করো.
\chThreeAnswer $\overline Y=\overline{A+\overline C}+\overline{\overline B+D}=\overline AC+B\overline D$। এখানে বাইরের AND বদলে OR হয়েছে।

\end{minipage}\par\smallskip
\Needspace{12\baselineskip}
\subsubsection*{ধরন ৪: সাধারণ গুণক নিয়ে পূরক জোড়া করা}
একই গুণক বের করো। অবশিষ্ট অংশে $X+\overline X=1$ বা $X\overline X=0$ থাকলে তা ব্যবহার করো।
\par\smallskip\noindent\begin{minipage}{\linewidth}
\chLogicExample{BA04.1}{৪.১}{এক চলক বাদ পড়া}
$F=\overline AB\overline C+\overline ABC+AD$ সরল করো.
\chThreeAnswer প্রথম দুই পদের $\overline AB$ এক; C শুধু দুটি সম্ভাবনাই প্রকাশ করছে।
\[\begin{aligned}F&=\overline AB\overline C+\overline ABC+AD\\
&=\overline AB(\overline C+C)+AD\\
&=\overline AB+AD\end{aligned}\]

\end{minipage}\par\smallskip
\par\smallskip\noindent\begin{minipage}{\linewidth}
\chLogicExample{BA04.2}{৪.২}{একাধিক পূরক জোড়া}
$F=A\overline B\overline C+A\overline BC+AB\overline C+ABC$ সরল করো.
\chThreeAnswer সব পদের A এক। B ও C-এর সব সম্ভাবনা থাকায় তাদের মান output নির্ধারণ করে না।
\[\begin{aligned}F&=A\overline B\overline C+A\overline BC+AB\overline C+ABC\\
&=A\overline B(\overline C+C)+AB(\overline C+C)\\
&=A(\overline B+B)\\
&=A\end{aligned}\]

\end{minipage}\par\smallskip
\Needspace{12\baselineskip}
\subsubsection*{ধরন ৫: বড় শর্তের মধ্যে ছোট শর্ত শোষিত হওয়া}
$X+XY=X$ এবং $X(X+Y)=X$। X একটি যৌগিক পদও হতে পারে।
\par\smallskip\noindent\begin{minipage}{\linewidth}
\chLogicExample{BA05.1}{৫.১}{যৌগিক পদের absorption}
$F=AC+ABC+ACD$ সরল করো.
\chThreeAnswer $1+B+D=1$। তাই বাড়তি অনুমতি B,D ফল বদলায় না।
\[\begin{aligned}F&=AC+ABC+ACD\\
&=AC(1+B+D)\\
&=AC\end{aligned}\]

\end{minipage}\par\smallskip
\par\smallskip\noindent\begin{minipage}{\linewidth}
\chLogicExample{BA05.2}{৫.২}{POS-এ absorption}
$F=(A+D)(A+D+B)(A+D+C)$ সরল করো.
\chThreeAnswer $X=A+D$ ধরলে $X(X+B)(X+C)=X$।
\[\begin{aligned}F&=(A+D)(A+D+B)(A+D+C)\\
&=(A+D)\end{aligned}\]

\end{minipage}\par\smallskip
\Needspace{12\baselineskip}
\subsubsection*{ধরন ৬: পূরকসহ elimination}
$X+\overline XY=X+Y$ এবং $X(\overline X+Y)=XY$। সরাসরি absorption-এর সঙ্গে এই সূত্রের পার্থক্য আছে।
\par\smallskip\noindent\begin{minipage}{\linewidth}
\chLogicExample{BA06.1}{৬.১}{যৌগিক Y নিয়ে elimination}
$F=\overline D+DAB$ সরল করো.
\chThreeAnswer $X=\overline D$, $Y=AB$। বণ্টন সূত্রে প্রথম গুণক ১ হয়।
\[\begin{aligned}F&=\overline D+DAB\\
&=(\overline D+D)(\overline D+AB)\\
&=\overline D+AB\end{aligned}\]

\end{minipage}\par\smallskip
\par\smallskip\noindent\begin{minipage}{\linewidth}
\chLogicExample{BA06.2}{৬.২}{AND শর্তে বিপরীত চলক}
$F=(A+C)(\overline A\,\overline C+B)$ সরল করো.
\chThreeAnswer $\overline{A+C}=\overline A\,\overline C$। X ও তার পূরকের গুণফল শূন্য।
\[\begin{aligned}F&=(A+C)(\overline A\,\overline C+B)\\
&=(A+C)\overline A\,\overline C+(A+C)B\\
&=B(A+C)\end{aligned}\]

\end{minipage}\par\smallskip
\Needspace{12\baselineskip}
\subsubsection*{ধরন ৭: POS প্রসারণ ও SOP-তে factor করা}
$X+YZ=(X+Y)(X+Z)$ এবং $(X+Y)(X+Z)=X+YZ$। সাধারণ বীজগণিতে প্রথম সূত্রটি একইভাবে চলে না।
\par\smallskip\noindent\begin{minipage}{\linewidth}
\chLogicExample{BA07.1}{৭.১}{সাধারণ OR পদ}
$F=(\overline A+B)(\overline A+C)$-কে SOP করো.
\chThreeAnswer $\overline A\,\overline A=\overline A$; তার সঙ্গে থাকা B,C পদ absorption-এ বাদ পড়ে।
\[\begin{aligned}F&=(\overline A+B)(\overline A+C)\\
&=\overline A+\overline AC+\overline AB+BC\\
&=\overline A+BC\end{aligned}\]

\end{minipage}\par\smallskip
\par\smallskip\noindent\begin{minipage}{\linewidth}
\chLogicExample{BA07.2}{৭.২}{একাধিক গুণকের factor}
$F=AB+CD$-কে POS-এ লেখো.
\chThreeAnswer দুইবার বণ্টন সূত্র প্রয়োগ করো। এই রূপের সব গুণক রাখতে হবে।
\[\begin{aligned}F&=AB+CD\\
&=(AB+C)(AB+D)\\
&=(A+C)(B+C)(A+D)(B+D)\end{aligned}\]

\end{minipage}\par\smallskip
\Needspace{12\baselineskip}
\subsubsection*{ধরন ৮: একাধিক স্তরের পূরক}
সবচেয়ে বাইরের পূরক খুলে এক স্তর করে এগোও। প্রতিটি overbar কোন পদের ওপর আছে তা চিহ্নিত করো।
\par\smallskip\noindent\begin{minipage}{\linewidth}
\chLogicExample{BA08.1}{৮.১}{ভেতরের পূরক উঠে যাওয়া}
$F=\overline{\overline{AB}+C}$ সরল করো.
\chThreeAnswer ভেতরের AB দ্বৈত পূরক হয়ে ফিরে আসে।
\[\begin{aligned}F&=\overline{\overline{AB}+C}\\
&=\overline{\overline{AB}}\,\overline C\\
&=AB\overline C\end{aligned}\]

\end{minipage}\par\smallskip
\par\smallskip\noindent\begin{minipage}{\linewidth}
\chLogicExample{BA08.2}{৮.২}{বন্ধনীর ওপর দ্বিতীয় পূরক}
$F=\overline{(A+\overline B)\overline{C+D}}$ সরল করো.
\chThreeAnswer বাইরের AND-এর পূরক OR; দ্বিতীয় পদের দ্বৈত পূরক উঠে যায়।
\[\begin{aligned}F&=\overline{(A+\overline B)\overline{C+D}}\\
&=\overline{A+\overline B}+\overline{\overline{C+D}}\\
&=\overline AB+C+D\end{aligned}\]

\end{minipage}\par\smallskip
\Needspace{12\baselineskip}
\subsubsection*{ধরন ৯: ধ্রুবক ০ বা ১ পাওয়া}
পরস্পরবিরোধী প্রয়োজন একসঙ্গে AND হলে output ০। সব সম্ভাবনাকে OR করলে output ১।
\par\smallskip\noindent\begin{minipage}{\linewidth}
\chLogicExample{BA09.1}{৯.১}{দুটি বিপরীত শর্ত}
$F=(A+C)\overline A\,\overline C$ সরল করো.
\chThreeAnswer একই সঙ্গে A অথবা C-কে ১ এবং দুটিকেই ০ চাওয়া হয়েছে।
\[\begin{aligned}F&=(A+C)\overline A\,\overline C\\
&=A\overline A\,\overline C+C\overline A\,\overline C\\
&=0\end{aligned}\]

\end{minipage}\par\smallskip
\par\smallskip\noindent\begin{minipage}{\linewidth}
\chLogicExample{BA09.2}{৯.২}{সব সম্ভাবনায় output ১}
$F=AB+A\overline B+\overline AC+\overline A\,\overline C$ সরল করো.
\chThreeAnswer A=১ ও A=০—দুটির জন্যই অন্য চলকের সব সম্ভাবনা রাখা হয়েছে।
\[\begin{aligned}F&=AB+A\overline B+\overline AC+\overline A\,\overline C\\
&=A(B+\overline B)+\overline A(C+\overline C)\\
&=A+\overline A\\
&=1\end{aligned}\]

\end{minipage}\par\smallskip
\Needspace{12\baselineskip}
\subsubsection*{ধরন ১০: Consensus-এ অপ্রয়োজনীয় পদ বাদ দেওয়া}
এক চলক বিপরীত অবস্থায় দুই পদে থাকলে তাদের অন্য অংশের গুণফল অনেক সময় আগেই অন্তর্ভুক্ত থাকে। পদ বাদ দেওয়ার আগে সূত্রের তিন অংশ মেলাও।
\par\smallskip\noindent\begin{minipage}{\linewidth}
\chLogicExample{BA10.1}{১০.১}{SOP consensus}
$F=AC+\overline A D+CD$ সরল করো.
\chThreeAnswer $CD=CD(A+\overline A)=ACD+\overline ACD$; নতুন পদদুটি প্রথম দুই পদের মধ্যে আছে।
\[\begin{aligned}F&=AC+\overline A D+CD\\
&=AC+\overline A D+ACD+\overline ACD\\
&=AC+\overline A D\end{aligned}\]

\end{minipage}\par\smallskip
\par\smallskip\noindent\begin{minipage}{\linewidth}
\chLogicExample{BA10.2}{১০.২}{POS consensus}
$F=(A+C)(\overline A+D)(C+D)$ সরল করো.
\chThreeAnswer প্রথম দুই বন্ধনী ১ হলে তৃতীয়টিও ১: A=০ হলে C=১ লাগবে; A=১ হলে D=১ লাগবে।
\[\begin{aligned}F&=(A+C)(\overline A+D)(C+D)\\
&=(A+C)(\overline A+D)\end{aligned}\]

\end{minipage}\par\smallskip
\Needspace{12\baselineskip}
\subsubsection*{ধরন ১১: রাশিতে XOR বা XNOR চেনা}
দুই চলকের বিপরীত মানের পদদুটি XOR; একই মানের পদদুটি XNOR। Gate-এর bubble থাকলে শেষ পূরকটিও দেখো।
\par\smallskip\noindent\begin{minipage}{\linewidth}
\chLogicExample{BA11.1}{১১.১}{POS থেকে XOR}
$F=(A+D)(\overline A+\overline D)$ সরল করে gate চিহ্নিত করো.
\chThreeAnswer $A\overline D+\overline AD=A\oplus D$।
\[\begin{aligned}F&=(A+D)(\overline A+\overline D)\\
&=A\overline A+A\overline D+D\overline A+D\overline D\\
&=A\overline D+\overline AD\end{aligned}\]

\end{minipage}\par\smallskip
\par\smallskip\noindent\begin{minipage}{\linewidth}
\chLogicExample{BA11.2}{১১.২}{XOR-এর পূরকে সমতা}
$F=\overline{B\overline C+\overline BC}$ সরল করো.
\chThreeAnswer এটি $\overline{B\oplus C}$ বা XNOR।
\[\begin{aligned}F&=\overline{B\overline C+\overline BC}\\
&=(\overline B+C)(B+\overline C)\\
&=\overline B\,\overline C+BC\end{aligned}\]

\end{minipage}\par\smallskip
\Needspace{12\baselineskip}
\subsubsection*{ধরন ১২: সূত্র প্রমাণ ও ভুল সমতা খণ্ডন}
সমতুল্যতা প্রমাণে সব input-এ output মিলতে হবে। একটি বিপরীত ফল পেলেই দাবিটি ভুল প্রমাণিত হয়।
\par\smallskip\noindent\begin{minipage}{\linewidth}
\chLogicExample{BA12.1}{১২.১}{সত্যক সারণিতে যৌগিক সূত্র}
$AB+A\overline B=A$ সত্যক সারণিতে দেখাও.
\chThreeAnswer প্রতিটি সারিতে শেষ দুই column সমান।
\begin{center}\begin{tabular}{|c|c|c|c|c|c|}\hline
$A$ & $B$ & $AB$ & $A\overline B$ & $AB+A\overline B$ & $A$\\\hline
0 & 0 & 0 & 0 & 0 & 0\\\hline
0 & 1 & 0 & 0 & 0 & 0\\\hline
1 & 0 & 0 & 1 & 1 & 1\\\hline
1 & 1 & 1 & 0 & 1 & 1\\\hline
\end{tabular}\end{center}

\end{minipage}\par\smallskip
\par\smallskip\noindent\begin{minipage}{\linewidth}
\chLogicExample{BA12.2}{১২.২}{একটি input দিয়ে ভুল ধরা}
এক শিক্ষার্থী লিখেছে $\overline{A+B}=\overline A+\overline B$। দাবিটি পরীক্ষা করো.
\chThreeAnswer $A=0,B=1$-এ বামপক্ষ ০, ডানপক্ষ $1+0=1$। তাই সমতা ভুল। সঠিক সূত্র $\overline{A+B}=\overline A\,\overline B$। শুধু 00 input পরীক্ষা করলে ভুলটি ধরা পড়ত না।

\end{minipage}\par\smallskip
\smallhead{ক--খ প্রশ্নের ধারণা}
বুলিয়ান $1+1=1$ কারণ অন্তত একটি সত্য শর্তে OR সত্য হয়। ডি-মরগ্যানের সূত্র দিয়ে বড় পূরক ভেঙে পৃথক input-এর পূরক এবং AND/OR দিয়ে সমতুল্য রাশি পাওয়া যায়। সরলীকরণে সত্যক সারণি অপরিবর্তিত রেখে প্রয়োজনীয় gate ও সংযোগ কমতে পারে। দ্বৈত রূপ এবং পূরক পৃথক প্রক্রিয়া।
